\documentclass[12pt]{article}
\usepackage[brazil]{babel}
\usepackage[utf8]{inputenc}
\usepackage[T1]{fontenc}
\usepackage{amsmath, amssymb}
\usepackage{geometry}
\usepackage{graphicx}
\usepackage{xcolor}
\usepackage{enumitem}

\definecolor{maincolor}{RGB}{0,100,50}
\usepackage{fancyhdr}
\definecolor{forestgreen}{RGB}{0,100,50}
\usepackage[most]{tcolorbox}

\newtcolorbox{block}[1]{%
  colback=green!5,
  colframe=maincolor,
  coltitle=white,
  title=#1,
  fonttitle=\bfseries,
  boxrule=1.5pt,
  arc=3mm,
  left=6pt,
  right=6pt,
  top=6pt,
  bottom=6pt
}

\geometry{margin=2.5cm}
\pagestyle{fancy}
\fancyhf{}

\renewcommand{\headrulewidth}{2pt}
\renewcommand{\footrulewidth}{2pt}

\renewcommand{\headrule}{\hbox to\headwidth{%
  \color{forestgreen}\leaders\hrule height \headrulewidth\hfill}}

\renewcommand{\footrule}{\hbox to\headwidth{%
  \color{forestgreen}\leaders\hrule height \footrulewidth\hfill}}

\setlength{\headheight}{32pt}

\fancyhead[L]{%
  % Certifique-se de que o arquivo de imagem esteja na mesma pasta
  \includegraphics[height=1.2cm]{Universidade_de_Rio_Verde.jpg}
}

\fancyhead[R]{%
  \raggedleft
  \textcolor{forestgreen}{%
    \small
    \textbf{Gabarito da Lista 3.}\\
    Me.~Sergio Souza Novak\\
    \textit{Lógica para Engenharia de Software}
  }
}

\fancyfoot[C]{%
  \textcolor{forestgreen}{\thepage}
}

\begin{document}

\vspace*{0.2cm}
\section*{Gabarito — Relações Binárias, Propriedades e Equivalência}

\begin{enumerate}

\item $A = \{1, 2, 3\}$ e $B = \{a, b\}$.
\begin{enumerate}[label=\alph*)]
    \item $A \times B = \{(1, a), (1, b), (2, a), (2, b), (3, a), (3, b)\}$.
    \item $B \times A = \{(a, 1), (a, 2), (a, 3), (b, 1), (b, 2), (b, 3)\}$.
    \item \textbf{Não são iguais}. Por exemplo, o par ordenado $(1, a) \in A \times B$, mas $(1, a) \notin B \times A$, visto que $1 \notin B$ e $a \notin A$. A ordem dos elementos em um par ordenado é fundamental.
\end{enumerate}

\item $A = \{1, 2, 3, 4, 6\}$ e $R$: ``$x$ divide $y$''.
\begin{enumerate}[label=\alph*)]
    \item $R = \{(1,1), (1,2), (1,3), (1,4), (1,6), (2,2), (2,4), (2,6), (3,3), (3,6), (4,4), (6,6)\}$.
    \item $M_R = \begin{pmatrix}
    1 & 1 & 1 & 1 & 1 \\
    0 & 1 & 0 & 1 & 1 \\
    0 & 0 & 1 & 0 & 1 \\
    0 & 0 & 0 & 1 & 0 \\
    0 & 0 & 0 & 0 & 1
    \end{pmatrix}$. \quad (\textit{As linhas e colunas representam os elementos 1, 2, 3, 4, 6 na ordem crescente}).
\end{enumerate}

\item Para $A = \{1, 2, 3, 4\}$:
\begin{enumerate}[label=\alph*)]
    \item $R_1 = \{(1,1), (2,2), (3,3), (4,4), (1,2), (2,1)\}$.\\
    Reflexiva (sim, contém todos os pares $(x,x)$), simétrica (sim, contém reflexos), antissimétrica (não, pois contém $(1,2)$ e $(2,1)$ com $1 \ne 2$), transitiva (sim). \emph{Irreflexiva e assimétrica também falham.}
    \item $R_2 = \{(1,2), (2,3), (1,3)\}$.\\
    Reflexiva (não), irreflexiva (sim), simétrica (não), antissimétrica (sim), assimétrica (sim), transitiva (sim, pois $(1,2),(2,3) \in R$ e $(1,3) \in R$).
    \item $R_3 = \{(1,1), (1,2), (2,1), (2,2), (3,4), (4,3)\}$.\\
    Reflexiva (não, faltam $(3,3)$ e $(4,4)$), simétrica (sim), antissimétrica (não), transitiva (não, pois contém $(3,4)$ e $(4,3)$, mas falta $(3,3)$).
    \item $R_4 = \varnothing$.\\
    Reflexiva (não), irreflexiva (sim, vacuidade), simétrica (sim, vacuidade), antissimétrica e assimétrica (sim, vacuidade), transitiva (sim, vacuidade).
\end{enumerate}

\item Uma relação $R$ é de equivalência sobre um conjunto $A$ se for simultaneamente:\\
\textbf{1. Reflexiva:} Para todo $x \in A$, $x\,R\,x$.\\
\textbf{2. Simétrica:} Para quaisquer $x, y \in A$, se $x\,R\,y$, então $y\,R\,x$.\\
\textbf{3. Transitiva:} Para quaisquer $x, y, z \in A$, se $x\,R\,y$ e $y\,R\,z$, então $x\,R\,z$.

\item Em $\mathbb{Z}$:
\begin{enumerate}[label=\alph*)]
    \item \textbf{É de equivalência} (a igualdade é o exemplo trivial de equivalência).
    \item \textbf{Não é de equivalência}, pois falha na simetria (ex: $1 \le 2$, mas $2 \not\le 1$).
    \item \textbf{É de equivalência} (reflexiva: $a$ tem mesma paridade que $a$; simétrica: se $a$ e $b$ têm mesma, $b$ e $a$ também; transitiva: mesmíssima lógica).
\end{enumerate}

\item Módulo 4 em $\mathbb{Z}$:
\begin{enumerate}[label=\alph*)]
    \item $[0] = \{ \dots, -4, 0, 4, 8, \dots \}$,\\
    $[1] = \{ \dots, -3, 1, 5, 9, \dots \}$,\\
    $[2] = \{ \dots, -2, 2, 6, 10, \dots \}$,\\
    $[3] = \{ \dots, -1, 3, 7, 11, \dots \}$.
    \item Os quatro subconjuntos são disjuntos entre si, não são vazios e sua união é igual a todo o conjunto de $\mathbb{Z}$ (pelo algoritmo de divisão euclidiana, qualquer inteiro deixa um resto entre 0 e 3).
\end{enumerate}

\item Partição $\mathcal{P} = \{\{a, b\}, \{c\}, \{d, e\}\}$ de $A = \{a, b, c, d, e\}$:
\begin{enumerate}[label=\alph*)]
    \item \textbf{Não-vazio:} nenhum conjunto é $\varnothing$. \textbf{Disjuntos:} interseções duas a duas dão $\varnothing$. \textbf{Cobertura união:} $\{a,b\} \cup \{c\} \cup \{d,e\} = A$. É partição.
    \item $R_\mathcal{P} = \{(a,a), (b,b), (a,b), (b,a), (c,c), (d,d), (e,e), (d,e), (e,d)\}$.
\end{enumerate}

\item Aplicação de Testes (Eng. Software):
\begin{enumerate}[label=\alph*)]
    \item Existem exatas \textbf{três} classes de equivalência (Criança, Adolescente, Adulto).
    \item O engenheiro precisa criar apenas \textbf{3 casos de teste} (um representante de cada classe, ex: 5 anos, 15 anos e 30 anos), pois o particionamento baseia-se na teoria matemática de que qualquer valor da mesma classe vai disparar os mesmos desvios condicionais e obter os mesmos comportamentos do sistema.
\end{enumerate}

\end{enumerate}

\end{document}
